\documentclass[12pt,a4paper]{article}
\usepackage[utf8]{inputenc}
\usepackage[T1]{fontenc}
\usepackage[italian]{babel}
\usepackage{libertinus}
\usepackage{amssymb}
\usepackage[a4paper,top=25mm,bottom=25mm,left=22mm,right=22mm,headheight=15pt]{geometry}
\usepackage{graphicx}
\usepackage[table]{xcolor}
\usepackage{tabularx,array}
\usepackage{titlesec}
\usepackage{titletoc}
\usepackage{microtype}
\usepackage{fancyhdr}
\usepackage{enumitem}
\usepackage[hidelinks]{hyperref}

% Numero del verbale: da impostare una sola volta (usato in intestazione, copertina e ID)
\newcommand{\vnum}{5}

% Palette colori del gruppo Bending Forks
\definecolor{bfblue}{HTML}{1B3F9E}
\definecolor{bfrow}{HTML}{F2F5FC}
\definecolor{bfline}{HTML}{C9D1E6}
\definecolor{bfhead}{HTML}{2E6BE6}

\newcommand{\tightlist}{\setlength{\itemsep}{1pt}\setlength{\parskip}{0pt}}
\setlist{topsep=2pt}
\setlist[itemize]{leftmargin=1.3em,label=\textbullet}
\setlist[enumerate]{leftmargin=1.6em}
\setlength{\parindent}{0pt}
\sloppy
\setlength{\parskip}{5pt}
\renewcommand{\arraystretch}{1.45}
\setcounter{secnumdepth}{3}
\setcounter{tocdepth}{3}

% Formattazione titoli sezioni
\titleformat{\section}{\LARGE\bfseries\color{bfblue}}{\thesection.}{0.3em}{}[{\color{bfblue}\vspace{-2pt}\titlerule[0.9pt]}]
\titleformat{\subsection}{\Large\bfseries\color{bfblue}}{\thesubsection.}{0.3em}{}
\titleformat{\subsubsection}{\normalsize\bfseries}{\thesubsubsection.}{0.3em}{}
\titlespacing*{\section}{0pt}{14pt}{6pt}
\titlespacing*{\subsection}{0pt}{10pt}{4pt}
\titlespacing*{\subsubsection}{0pt}{10pt}{3pt}

% Indice cliccabile con puntini/trattini
\contentsmargin{2.2em}
\newcommand{\bfleader}{\titlerule*[0.32em]{\textendash}}
\titlecontents{section}[0pt]{\addvspace{5pt}\large\bfseries}{\thecontentslabel.\ }{}{\hspace{0.4em}\bfleader\contentspage}
\titlecontents{subsection}[1.2em]{\addvspace{0pt}}{\thecontentslabel.\ }{}{\hspace{0.4em}\bfleader\contentspage}
\titlecontents{subsubsection}[2.4em]{\addvspace{0pt}}{\thecontentslabel.\ }{}{\hspace{0.4em}\bfleader\contentspage}

% Intestazioni e piè di pagina
\pagestyle{fancy}
\fancyhf{}
\fancyhead[L]{\footnotesize\sffamily\bfseries\color{bfhead}\textls[120]{BENDING FORKS}}
\fancyhead[R]{\footnotesize\sffamily Verbale Interno N. \vnum}
\fancyfoot[C]{\small Pagina \thepage}
\renewcommand{\headrulewidth}{0.9pt}

% Colorazione alternata per le tabelle
\newcommand{\bftable}{\rowcolors{2}{bfrow}{white}}
\usepackage{etoolbox}
\AtBeginEnvironment{tabularx}{\bftable}

\begin{document}

% ---------- Copertina ----------
\begin{titlepage}
\centering
\vspace*{4mm}
\includegraphics[width=105mm]{Assets/logo-bending-forks-trasparente.png}\\[4mm]
\underline{bending.forks.swe@gmail.com}

\vspace{42mm}
{\fontsize{24}{28}\selectfont\bfseries Verbale Interno N. \vnum\par}
\vspace{4mm}
{\Large\color{bfblue} Riunione del 28/10/2026\par}

\vfill
\begin{tabular}[t]{r|l}
\textbf{Responsabile} & Aldo Toffoli \\
\textbf{Scriba / Redattore} & Luca Stangherlin \\
\textbf{Moderatore} & Aldo Toffoli \\
\textbf{Verificatore} & Filippo Rampazzo \\
\end{tabular}
\hspace{14mm}
\begin{tabular}[t]{r|l}
\textbf{Versione} & 1.0 \\
\textbf{Tipo} & Verbale Interno \\
\textbf{Destinatari} & \emph{Gruppo Bending Forks} \\
 & \emph{Prof. Tullio Vardanega} \\
 & \emph{Prof. Riccardo Cardin} \\
\end{tabular}

\vspace{12mm}
{\large\emph{\textbf{Abstract dei contenuti:}}}\\[1mm]
\begin{minipage}{150mm}\centering Sintesi e decisioni operative emerse durante la riunione interna del team del giorno 28/10/2026, riguardante la pianificazione delle attività e l'avanzamento del progetto.\end{minipage}
\vspace*{6mm}
\end{titlepage}

% ---------- Indice ----------
\vspace*{2mm}
{\fontsize{22}{26}\selectfont\bfseries Indice\par}
\vspace{6mm}
\makeatletter\@starttoc{toc}\makeatother
\newpage

% ---------- Registro delle versioni ----------
\hypertarget{registro-delle-versioni}{%
\section{Registro delle versioni}}

\begin{center}
\begin{tabularx}{\linewidth}{>{\hsize=0.60\hsize\raggedright\arraybackslash}X>{\hsize=0.75\hsize\raggedright\arraybackslash}X>{\hsize=0.95\hsize\raggedright\arraybackslash}X>{\hsize=1.70\hsize\raggedright\arraybackslash}X}
\rowcolor{bfblue}\textcolor{white}{\textbf{Versione}} & \textcolor{white}{\textbf{Data}} & \textcolor{white}{\textbf{Autore}} & \textcolor{white}{\textbf{Descrizione}} \\
1.0 & 28/10/2026 & Luca Stangherlin & Creazione e redazione iniziale del documento \\ \arrayrulecolor{bfline}\hline
1.0 & 30/10/2026 & Filippo Rampazzo & Verifica e approvazione \\ \arrayrulecolor{bfline}\hline
\end{tabularx}
\end{center}

% ---------- Informazioni generali ----------
\hypertarget{informazioni-generali}{%
\section{Informazioni generali}}

\hypertarget{durata-e-luogo}{%
\subsection{Durata e luogo}}

\begin{itemize}
\tightlist
\item \textbf{Data:} 2026-10-28
\item \textbf{Inizio:} 15:15
\item \textbf{Fine:} 16:30
\item \textbf{Luogo / Modalità:} in remoto, tramite Discord
\end{itemize}

\hypertarget{partecipanti}{%
\subsection{Partecipanti}}

\begin{center}
\begin{tabularx}{\linewidth}{>{\raggedright\arraybackslash}X>{\centering\arraybackslash}X>{\centering\arraybackslash}p{2.2cm}>{\centering\arraybackslash}p{2.2cm}}
\rowcolor{bfblue}\textcolor{white}{\textbf{Nome e Cognome}} & \textcolor{white}{\textbf{Ruolo Progetto}} & \textcolor{white}{\textbf{Presente}} & \textcolor{white}{\textbf{Assente}} \\
Marco Beltrame & Membro & $\blacksquare$ & $\square$ \\ \arrayrulecolor{bfline}\hline
Lorenzo Milanese & Membro & $\blacksquare$ & $\square$ \\ \arrayrulecolor{bfline}\hline
Filippo Rampazzo & Verificatore & $\blacksquare$ & $\square$ \\ \arrayrulecolor{bfline}\hline
Nicholas Salmaso & Membro & $\blacksquare$ & $\square$ \\ \arrayrulecolor{bfline}\hline
Luca Stangherlin & Membro & $\blacksquare$ & $\square$ \\ \arrayrulecolor{bfline}\hline
Aldo Toffoli & Responsabile & $\blacksquare$ & $\square$ \\ \arrayrulecolor{bfline}\hline
Matteo Wu & Membro & $\square$ & $\blacksquare$ \\ \arrayrulecolor{bfline}\hline
\end{tabularx}
\end{center}

% ---------- Ordine del giorno ----------
\hypertarget{ordine-del-giorno}{%
\section{Ordine del giorno}}

\begin{enumerate}
\tightlist
\item Revisione e completamento del documento \emph{Valutazione dei capitolati}
\item Piano di massima, rischi principali e mitigazioni
\item Candidatura: lettera di presentazione e preventivo
\item Sito web del gruppo
\end{enumerate}

% ---------- Discussione e decisioni ----------

\newpage
\hypertarget{discussione-e-decisioni}{%
\section{Discussione e decisioni}}

\hypertarget{argomento-1}{%
\subsection{Revisione del documento Valutazione dei capitolati}}

\textbf{Discussione:} \\
Il gruppo ha rivisto il documento di valutazione dei capitolati. Sono stati collettivamente asseganti dei  voti (da 1 a 5, con media dei pareri) su: interesse tecnico, complessità, chiarezza dei requisiti, qualità del rapporto con la proponente, fattibilità tecnica e allineamento delle competenze. 
Il documento è stato infine revisionato dai componenti del team. 

\textbf{Decisione presa:}
\begin{itemize}
\tightlist
\item[$\checkmark$]  I voti della tabella di confronto sono assegnati collettivamente, su una scala da 1 a 5, coerentemente con i giudizi espressi nella \textit{Valutazione dei Capitolati}.

\item[$\checkmark$]  Si mantiene la lunghezza attuale del documento (circa dieci pagine per dieci capitolati).
\end{itemize}

\textbf{Motivazione:} \\
Coerenza interna tra tabella, descrizioni e conclusioni.

\textbf{Responsabili attuazione:} Aldo Toffoli (redazione); Marco Beltrame (revisione) \\
\textbf{Scadenza:} 2026-10-28


\hypertarget{argomento-4}{%
\subsection{Candidatura: lettera di presentazione e preventivo}}

\textbf{Discussione:} \\
 La lettera di presentazione e il preventivo sono stati stilati in sotto-gruppi; la revisione di entrambi viene svolta collettivamente da tutti partecipanti alla riunione con il seguente esito: la lettera viene approvata in attesa del preventivo; il preventivo è curato da Luca Stangherlin e Nicholas Salmaso che abbozzano un modello con paramentri e sezioni da definire con il resto del team. 

\textbf{Decisione presa:}
\begin{itemize}
\tightlist
\item[$\checkmark$]  Lettera di presentazione e preventivo sono completati entro il 29/10/2026. La lettera è datata al giorno del caricamento sul repository.
\item[$\checkmark$]  Il preventivo verrà completato dal gruppo il 29/10, partendo dallo scheletro su Drive e definendo collettivamente le ore per ruolo.
\end{itemize}

\textbf{Motivazione:} \\
Controllare la conformità rispetto ai canoni di ogni componente del team dei docuementi per la candidatura. 

\textbf{Responsabili attuazione:} Lorenzo Milanese e Filippo Rampazzo (lettera); Luca Stangherlin e Nicholas Salmaso (preventivo) \\
\textbf{Scadenza:} 2026-10-29

\hypertarget{argomento-5}{%
\subsection{Sito web del gruppo}}

\textbf{Discussione:} \\
Viene definita l'esigenza di realizzare un sito web per la consultazione dei documenti in maniera più rapida e intuitiva. La priorità è stabilita media.

\textbf{Decisione presa:}
\begin{itemize}
\tightlist
\item[$\checkmark$]  Il sito è realizzato con GitHub Pages e ha priorità minore di lettera e preventivo.
\end{itemize}

\textbf{Motivazione:} \\
 La priorità media evita di sottrarre tempo agli elementi obbligatori della candidatura.

\textbf{Responsabili attuazione:} Lorenzo Milanese; Filippo Rampazzo \\
\textbf{Scadenza:} 2026-10-30

\hypertarget{varie-ed-eventuali}{%
\subsection{Varie ed eventuali}}
\begin{itemize}
\tightlist
\item \textbf{Versionamento:} Lorenzo Milanese aggiorna il verbale esterno in base alle rilevazioni di Vimar. La verifica è affidata a Luca Stangherlin.
\end{itemize}

% ---------- Azioni da intraprendere (TODO) ----------
\hypertarget{azioni-da-intraprendere}{
\section{Azioni da intraprendere (TODO)}}

\begin{center}
\begin{tabularx}{\linewidth}{>{\hsize=2.0\hsize\raggedright\arraybackslash}X >{\hsize=1.0\hsize\raggedright\arraybackslash}X >{\hsize=0.5\hsize\centering\arraybackslash}X >{\hsize=0.5\hsize\centering\arraybackslash}X}
\rowcolor{bfblue}\textcolor{white}{\textbf{Descrizione}} & \textcolor{white}{\textbf{Responsabile}} & \textcolor{white}{\textbf{Scadenza}} & \textcolor{white}{\textbf{Priorità}} \\
Completare il preventivo (ripartizione delle ore per ruolo) e la dichiarazione di impegni con tutto il gruppo & Luca Stangherlin, Nicholas Salmaso & 2026-10-29 & Alta \\ \arrayrulecolor{bfline}\hline
Completare la lettera di presentazione (costo, data finale, contatti), farla firmare e caricarla sul repository & Lorenzo Milanese, Filippo Rampazzo & 2026-10-29 & Alta \\ \arrayrulecolor{bfline}\hline
Correggere il documento di revisione dei capitolati secondo le indicazioni del gruppo & Aldo Toffoli, Marco Beltrame & 2026-10-29 & Alta \\ \arrayrulecolor{bfline}\hline
Redigere il verbale della riunione del 28/10/2026 & Luca Stangherlin & 2026-10-29 & Media \\ \arrayrulecolor{bfline}\hline
Verificare il documento v1.1 con le correzioni sui commenti di Vimar & Luca Stangherlin & 2026-10-29 & Media \\ \arrayrulecolor{bfline}\hline
Realizzare il sito su GitHub Pages in un repository dedicato & Lorenzo Milanese, Filippo Rampazzo & 2026-10-30 & Media\\ \arrayrulecolor{bfline}\hline
\end{tabularx}
\end{center}

% ---------- Prossima riunione ----------
\hypertarget{prossima-riunione}{%
\section{Pianificazione prossima riunione}}

\begin{itemize}
\tightlist
\item \textbf{Data proposta:} 2026-10-29
\item \textbf{Ora prevista:} 15:00
\item \textbf{Luogo / Canale:} Virtuale, canale Discord 
\item \textbf{OdG preliminare:} Completamento del preventivo; finalizzazione e caricamento della lettera di presentazione; test e finiture del sito
\end{itemize}

\vspace{35pt}
\noindent
\begin{minipage}[t]{0.45\textwidth}
\textbf{Firma Scriba:}\\[18pt]
\makebox[0pt][l]{\raisebox{0mm}{\includegraphics[width=40mm]{Assets/firma_Luca.png}}}\rule{55mm}{0.5pt}\\[2pt]
\small Luca Stangherlin
\end{minipage}
\hfill
\begin{minipage}[t]{0.45\textwidth}
\begin{flushright}
\textbf{Firma Verificatore:}\\[18pt]
\makebox[0pt][l]{\raisebox{0mm}{\includegraphics[width=40mm]{Assets/firma_fil.jpeg}}}
\rule{55mm}{0.5pt}\\[2pt]
\small Filippo Rampazzo
\end{flushright}
\end{minipage}

\end{document}